\scanpage{39}
\uline{Récapitulation partielle}

À un module $M$ libre de rang $2$ sur $\mathbb Z/4\mathbb Z$ on associe

\noindent\makebox[0pt][r]{(49)\quad}\(\left\{\begin{minipage}{.77\linewidth}
a) L'octaèdre (non orienté) \inserted{de sommets} $\Delta_M=M^*/\pm1\ \bullet$, l'ens des trois diagonales s'identifiant à $(M/2M)^*\simeq M_0^*=J$

\medskip b) L'ensemble à quatre éléments $I$, qu'on peut p.ex. déduire de $\Delta$ comme
\[
 I_M\simeq\underbrace{F(\Delta)}_{\substack{\text{faces de }\Delta\\\text{i.e. sections de}\\\Delta\text{ sur }J}}/\text{antipodisme}
\]
\end{minipage}\right.\)

\reversemarginpar\margin{On a un iso. can. $J\simeq\mathcal P_{22}(I_M)$, $I_M$ torseur sous $\underset{\simeq\mathbb F_2^J/\mathbb F_2}{\underline{V_J(\mathbb F_2)\simeq M_0}}$ $I_M\simeq P_M/\mathbb F_2^{\underline\omega}$}\normalmarginpar

\margin{\(\left|\begin{array}{ll}\text{\uline{NB}}&F(\Delta)\simeq I\times\underline C_\Delta\\\bullet&\hfill\substack{\downarrow\\\text{\textquotedblleft couleurs\textquotedblright}\\\text{de faces}}\end{array}\right.\)}

\noindent\makebox[0pt][r]{\rule[-1.4ex]{.4pt}{4.6ex}\quad}c) Le carré $Q_M$ \deleted{$\unsure{S_M|S}$} \inserted{ayant comme ens de sommets} $=P_M/M_0$, où $P_M$ est le $\mathrm{End}(M_0)$-torseur formé des sections de $M^*$ sur $M_0^*$ satisfaisant $\unsure{\xi_i+\xi_j+\xi_h|\xi_i+\xi_j+\xi_k}=0$, et où on identifie $\mathrm{End}(M_0)$ à $M_0\oplus\mathbb F_2^{\underline\omega}$, de sorte que $\mathbb P_M/M_0$ est un $\mathbb F_2^{\underline\omega}$ torseur, donc est l'ens des sommets d'un carré \inserted{$Q=Q_M$} dont l'ens des codiagonales est $\underline\omega$

\reversemarginpar\margin{où $\underline\omega\simeq\underline\omega_I\simeq\underline\omega_J$}\normalmarginpar

Les données a) et b) ne dépendent que de la \deleted{fibré} droite projective $X_M=\mathbb P(M)$ sur $\mathbb Z/4\mathbb Z$ ($\Delta$ est l'ens des p\textsuperscript{ts} de celle-ci, $J$ \unsure{celui|celle|c'lui|de l'ens|l'ens} des p\textsuperscript{ts} de $(X_M)_0$ droite projective sur $\mathbb F_2$), mais non pas $Q_M$. On a des bijections canoniques

\[
\left\{\begin{aligned}
\underbrace{C_Q}_{\text{codiag de }Q}&\simeq
\underbrace{\underline\omega_\Delta}_{\deleted{\substack{\text{\unsure{couleurs des}}\\\text{faces de }\Delta}}}
\overset{\text{déf}}{\simeq}\underline\omega_I\simeq\underline\omega_J
&&\text{soit }\underline\omega\\
\underset{\text{diag de }Q}{D_Q}&\simeq
\underset{\text{\textquotedblleft couleurs\textquotedblright\ des faces}}{\underline C_\Delta}
\simeq\underline\omega\wedge\underline\mu
&&\text{où }\underline\mu\ \text{\unsure{défini|définie} ci-dessous}\\
\underset{\text{or. de }Q}{\underline\omega_Q}&\simeq\underline{\mathrm{or}}(\Delta)
\simeq\det(M)^* &&\text{soit }\underline\mu
\end{aligned}\right. \tag*{(50)}
\]
% archive sheet number: 38
